\subsection{One finite endpoint and one infinite endpoint}
\label{rw:infinity:one-finite-section}

Label the infinite endpoint by $q\to\infty$ and the finite
outgoing input by $p\to p_*>0$.
If this labeling reverses the fixed-field direction, both signs of
the resulting
$\ell$ are retained on the same actual positive-central passing set;
reflection is transferred back by the inverse return germ of Section~2.
The lower restrictions are those of
Lemma~\ref{rw:finite:lower-cones}. The upper alternatives are a
strict real pair with $\bar m<0$, a repeated pair
$\bar m=-\gamma^2$, $\bar c=2\gamma>0$, or the weak boundary
$\bar m=0$, $\bar c>0$. In the repeated case $p_*\ge1/\gamma$
is retained, since the strict barrier does not exclude these larger
inputs.

At a regular finite lower endpoint, incoming contraction dominates
any polynomial growth of the upper derivative in the large exit.
The next lemma gives this polynomial bound through a repeated
eigenvalue. Lower losses will then require a comparison of the two
slow rates.

\begin{lemma}[Upper expansion with negative $\bar m$]
\label{rw:infinity:polynomial-upper}
On a compact coefficient block with $\bar m<-m_0<0$ and
$\bar c>c_0>0$, and a compact interval of positive finite inputs,
every actual upper return with sufficiently large exit $q$ satisfies
$0<q_p\le C(1+q)^N$. The constants are uniform for small scale
and either sign of $\nu$, including at zero discriminant.
\end{lemma}
\begin{proof}
Put $Z=Y+\epsilon$, $\theta=\epsilon\bar B$,
$d=\bar c+\nu$ and
\[
 g(Z)=dZ-\epsilon(\bar c+2\nu)+\nu\epsilon^2/Z,
 \quad U=\frac{Z^{1-\theta}}{1-\theta},\quad
 W=Z^{-\theta}(X+g(Z)).
\]
Divide time by $Z$. Differentiating these coordinates in the original
upper field gives
\[
 U'=W,\qquad W'=-G(U)+f(Z)W,
\]
where
\begin{align*}
 f(Z)&=(1-2\theta)d+
 \frac{\epsilon\bar a+2\theta\epsilon(\bar c+2\nu)}Z
       -\frac{(1+2\theta)\nu\epsilon^2}{Z^2},\\
 G(U)&=Z^{-\theta-1}Y
       [1+\epsilon\bar a e_Y-(\bar m+\theta e_Y^2)Y],
 \qquad e_Y=\bar c+\nu Y/Z.
\end{align*}
Since $|f|\le C$ and $G>0$ for $Y>0$, every zero of $W$ is a
downward crossing. A returning arc therefore has one height maximum.
For large $Z$, separate the linear part:
\[
 \binom{U}{W}'=\mathcal L\binom{U}{W}+\binom0{\mathcal E},\qquad
 \mathcal L=\begin{pmatrix}0&1\\-k&\tau\end{pmatrix},\quad
 k=-(1-\theta)(\bar m+\theta d^2)>0,\quad
 \tau=(1-2\theta)d>0,
\]
where $|\mathcal E|\le CZ^{-1}(|U|+|W|)$.
The eigenvalues of $\mathcal L$ have uniformly positive real parts,
since $k,\tau>0$ on this compact block. The Lyapunov integral
\[
 \mathcal S=\int_0^\infty e^{-\mathcal L^{\mathsf T}t}
                                  e^{-\mathcal Lt}\,dt
\]
gives a positive matrix with uniform bounds for $\mathcal S$ and
$\mathcal S^{-1}$, also at zero discriminant, and
$\mathcal L^{\mathsf T}\mathcal S+\mathcal S\mathcal L=I$.
Choose a fixed large
$Z_0$. Then
$\mathcal V=(U,W)^{\mathsf T}\mathcal S(U,W)$ satisfies
$\mathcal V'\ge c\mathcal V$ on $Z\ge Z_0$.

On each monotone-height portion below $Z_0$,
\[
 \left|\frac{d\sqrt{1+W^2}}{dU}\right|
 =\left|\frac{-G+fW}{\sqrt{1+W^2}}\right|\le C.
\]
The positive margin of the initial value
$W=\epsilon^{-\theta}p$ bounds the outgoing time in a small initial
strip.
Away from its bottom, $G$ has a positive margin: whenever positive
$W$ is small, $W'<-c$, so that such a visit has bounded duration
and cannot be repeated before the height maximum. Otherwise $W$
stays away from zero. On the incoming low strip the preceding
inequality gives $|W|\ge cq-C$, so its time is $O(1/q)$.
For large $q$ the arc must reach $Z_0$, since the same inequality
would bound $q$ if it stayed below it. At its first crossing
$\mathcal V\ge c$, and at its last crossing
$\mathcal V\le C(1+q^2)$. Exponential growth on the high part and
the low-part bounds yield $T\le C(1+\log(1+q))$.
The divergence of the $(U,W)$ equation is $f$. Its section determinant
formula gives
\[
 q_p=\frac pq\exp\int_0^T f\,dt\le C(1+q)^N.
\]
The endpoint factor $\epsilon^{-\theta}$ cancels exactly. The
quadratic estimate uses no inverse eigenvalue difference, which gives
the required uniformity at the repeated eigenvalue.
\end{proof}

\begin{lemma}[A coarse count for attracting collar cycles]
\label{rw:infinity:attracting-count}
Suppose every periodic orbit in the collar whose ray labels lie
between two retained label bounds has return multiplier less than
one. Require the same property for every periodic orbit nested
between two of them, without assuming isolation beforehand. Then
there are at most four such cycles, all hyperbolic attracting.
\end{lemma}
\begin{proof}
The family \eqref{rw:family} has at most four affine equilibria.
For $\ell\ne0$, elimination using
$x=-y(dy+c)/(1+y)$ gives $yR(y)=0$, where
\[
 R(y)=-(1+y)^2+By(dy+c)^2+my(1+y)^2-a(dy+c)(1+y).
\]
This has degree at most three and $R(0)=-(1+ac)\ne0$;
there is no equilibrium on $y=-1$, since $Q(x,-1)=\ell$.
The elimination divides only by $1+y$, so a zero of $dy+c$
elsewhere simply gives $x=0$ and is retained. For $\ell=0$, the branches
$y=-1$ and $x=-cy$ give respectively
$Bx^2+ax+1+m=0$ and
$y[(Bc^2+m)y-(1+ac)]=0$. Each has at most two roots, since
$1+m$ and $1+ac$ have positive margins. The origin lies inside
every counted cycle, leaving at most three equilibria between them.

Two nested collar cycles bound a compact affine annulus in the collar;
the connected southern hemisphere lies outside it. If the annulus has
no equilibrium, the boundary of the inner attracting cycle's basin
is a nonempty compact invariant subset of its interior, separated from
both attracting boundary collars. The Poincare--Bendixson theorem
\cite[Chapter~7]{Teschl2012} gives a periodic orbit on that boundary.
It cannot attract a neighborhood, since that would make its basin
boundary point interior to a disjoint basin. Its intermediate ray labels
contradict the hypothesis. Hence the disjoint annuli between any finite
ordered selection of cycles each contain an equilibrium. There are at
most three such equilibria and therefore at most four cycles, even
without an initial finiteness assumption.
\end{proof}

\begin{proposition}[A regular finite arc with negative $\bar m$]
\label{rw:infinity:negative-regular}
If $\bar m<0$, $\bar c>0$ and $F_+>0$ on $[0,p_*]$, then
near this one-infinite configuration all counted cycles are simple
and attracting and there are at most four of them. On each common
positive-central passing interval the ratio of incoming to outgoing
central derivatives is less than $1/2$, so $K<0$ and there is at most
one closing zero on that interval. The coarse bound of four does not
require connected saturation. The statement holds for either sign of $\ell$.
\end{proposition}
\begin{proof}
Choose a fixed incoming crosscut $R$ before its first root with
$A_-(R)>A_{\rm out}+3$. Such a choice is possible: the action
diverges at a first root, and if no root occurs it diverges as
$R\to\infty$, since $F_-(R)\le C(1+R+R^2)$.
Lemma~\ref{rw:finite:incoming-contraction} applies on the actual
common positive-central passing set. Its power
$[(1+R^2)/(1+q^2)]^{1/(24\epsilon)}$ absorbs the fixed polynomial
from Lemma~\ref{rw:infinity:polynomial-upper}. The ratio of the
complete incoming and outgoing central derivatives is therefore
less than $1/2$. At any periodic closing point this is the return
multiplier, including one not initially assumed isolated.
By the ray ordering of Lemma~\ref{rw:geometry:endpoints}, an
intermediate collar cycle has input in the same finite interval and
exit above the same large $q$ threshold. Thus every intermediate
cycle satisfies the multiplier estimate, and
Lemma~\ref{rw:infinity:attracting-count} gives the bound of four
across all passing components. On each common passing interval,
$K<0$ gives the stated one-zero bound directly.
\end{proof}

\subsection{A strict real upper pair and a finite lower loss}

When the finite lower arc loses its exit, its contraction can compete
with the upper passage. We compare their logarithmic rates, beginning
with $\bar m<0$, $\bar c>0$ and $\bar c^2+4\bar m>0$.
In the projective field
\eqref{rw:upper:projective} put
\[
 b=1-\epsilon\bar B,\quad d=\bar c+\nu,\quad
 \Delta=\sqrt{d^2+4b\bar m},\quad
 \gamma_s=\frac{d(1-2\epsilon\bar B)-\Delta}{2b},\quad
 r_s=\gamma_s/\Delta>0.
\]
At $R=0$, $\varphi=(-d-\Delta)/(2b)$ the eigenvalues are
$-\gamma_s,\Delta$. Its strong intersection is
$P_s=p_c+O_N(\epsilon(1+|\log\epsilon|))+\nu O_N(1)$,
and its returning side is $p<P_s$.
The analytic invariant axes, constructed on a small projective
rectangle with fixed hyperbolic margins and continued through the
compact arc and low layer, separate the returning quadrant from the
adjacent node. Axis normalization and the reciprocal integral give
the complete radial output
\[
 s(t)=k t^{r_s}e^{w_s(t)},\qquad k>c,\qquad
 |D^i\partial_{\mathbf b}^{\alpha}w_s|\le C_Nt^{\alpha_0}.
\]
Here $D=t\partial_t$ and a fixed $\alpha_0>0$ is smaller than
half the minimum of $1,r_s$ on the chosen block. Integrals of
$t^{r_s}s^{-r_s}$ and their finite logarithmic derivatives prove
this bound also when a higher saddle resonance is present; the
strict exponent reduction absorbs those logarithms. Fix all section
and coefficient derivative orders needed below before
choosing the rectangle.

If $\bar B>0$, the incoming half-line is regular whenever the
finite outgoing lower arc loses positivity, since then $\bar a<0$.
The horizontal preparation is
$D_I(s)=K_Is^{h_I/\epsilon}e^{P_I(s)/\epsilon}$,
$h_I=b/\bar B$, $P_I(0)=0$, with bounded selected jets of $P_I$.
The bounded amplitude is $P_I$; its quotient by $\epsilon$ need not
be bounded. Composing with the upper passage gives
\begin{equation}\label{rw:infinity:strict-rate}
 J(t)=K_Jt^{h_\infty/\epsilon}e^{\Psi(t)/\epsilon},\quad
 h_\infty=r_sh_I>c,\qquad
 \epsilon tJ''/J'=h_\infty-\epsilon+E(t),
\end{equation}
where $\Psi=h_Iw_s+P_I(s(t))$ after constants have been removed,
and
$E=D\Psi+\epsilon D^2\Psi/(h_\infty+D\Psi)$.
Its prescribed Euler and coefficient jets are $O(t^{\alpha_0})$.
The denominator has a positive margin on a fixed smaller section.

\begin{proposition}[Strict real pair with lower loss]
\label{rw:infinity:strict-lower-loss}
For $\bar B>0$, a first simple lower root at $p_c$ has local
closing count three. A double lower root at $p_c$ also has count
three; a double lower root strictly before $p_c$ has count two.
The bounds include multiplicity and either sign of $\ell$.
\end{proposition}
\begin{proof}
At a simple root choose its regular power input as
$u=\delta+t$, where $t>0$ is the upper returning coordinate.
The exact splitting is the difference of the two strong intersections
in this coordinate. Write $\lambda_+=h_+/\epsilon$ and
$g=h_+-h_\infty$. The physical interval has
$t>0$, $\delta+t>0$, and positive central outputs. The logarithmic
critical function has the full cleared numerator
\[
 M=\delta(h_\infty-\epsilon)-gt+(\delta+t)E(t)
                       -\epsilon t(\delta+t)R_+(\delta+t),
 \qquad M=t(\delta+t)K'.
\]
Here $R_+$ is the bounded ordinary correction to
$D_+''/D_+'=(\lambda_+-1)/u+R_+(u)$.

On the zero-scale matching surface, with $p=p_c$, the exponent
difference is
\[
 g_0=\frac{p^2(\bar B+\bar m)}
 {\bar B(1-\bar Bp^2)(1+\bar mp^2)}.
\]
Indeed the lower exponent is $p^2/(1-\bar Bp^2)$ and the
incoming exponent is $-\bar mp^2/[\bar B(1+\bar mp^2)]$.
At $\bar m=-\bar B$, $\bar a=-\bar c$, the Jacobian of
$(\delta_{\rm raw},g)$ in $(\bar a,\bar m)$ tends to
$p^4/[\bar B(1-\bar Bp^2)^3]>0$. This follows by multiplying
$\partial_{\bar a}\delta_{\rm raw}=p^2/(1-\bar Bp^2)$ by
the derivative of $g_0$ along the matching surface.
The input change multiplies this determinant by a positive unit.
Thus $(\delta,g,\nu)$ give uniform inverse coefficient boxes.
The second real center of Proposition~\ref{rw:real-centers}
lies exactly on $\delta=g=\nu=0$: its actual maps agree up to
the common strong endpoint, and comparison of their positive leading
powers forces $\lambda_+=r_sh_I/\epsilon$ there. Conversely the
nonzero Jacobian makes this the local center axis.

The center identity removes the constant transverse term in the
critical numerator. Extend the relative coefficients to the larger
box as in Section~5 and divide there to obtain
\[
 M=\delta A(t)-gtB(t)+\nu tQ(t),\quad A,B>c,
 \quad |A'|+|B'|\le Ct^{\alpha_0-1},
 \quad |Q|+|tQ'|\le C.
\]
For example first set $\delta=0$ and divide $M/t$ in $g,\nu$,
then integrate its $\delta$ derivative. The constant columns are
$h_\infty-\epsilon$ and $-1$, and the differentiated remainders
are the Euler-small terms of \eqref{rw:infinity:strict-rate}.
Each absolute lower map has a fixed positive action before its
central output. Its finite derivatives are exponentially small;
center division and Lemma~\ref{rw:centers:relative-forcing} therefore
give $|\nu|=o(1)(|\delta|+|g|)$ on a fibre with a closing zero.
At a zero of $M$ these estimates force
$\delta g>0$ and $|\delta|\asymp|g|t$; differentiating the
displayed division then gives sign $M'=-\operatorname{sign}g$.
All zeros are therefore simple crossings in the same direction, so
there is at most one. Away from $g_0=0$, the undivided numerator
gives the same conclusion after a fixed section restriction.
Since $M$ is a positive multiple of $K'$, two Rolle steps give
three closing zeros, with multiplicity.

For a double root, retain the complete moving slow input $u(t)$,
$u'>c>0$, and $f_L=\mu_L+u^2$. Its rate after the regular input
change has numerator
$\widehat A=A_L(u)u'+\epsilon f_Lu''/u'>c$.
The whole-path condition on $f_L$ is imposed. The exact critical
numerator is
\[
 M=f_L(h_\infty-\epsilon+E)-t\widehat A.
\]
At a double root at $p_c$, $u(0)$ tends to zero. At a zero,
$f_L\asymp t$, and differentiation gives
\[
 M'=2uu'(h_\infty-\epsilon+E)+f_LE'
                      -\widehat A-t\widehat A'<0.
\]
Indeed $tE'=O(t^{\alpha_0})$, and all other positive-size errors
are small compared with $\widehat A$. Hence $M$ has at most one
zero. If the double root is strictly before $p_c$, the limiting
entry is negative and $f_L(0)$ has a positive margin, so $M>0$
on a smaller section. Thus the critical numerator has at most one zero
at contact and none in the separated case, giving the asserted counts. The scalar and
central-output restrictions are monotone passing conditions, whose
intersection is one interval.
\end{proof}

\subsection{Repeated upper roots and signed horizontal parameters}

At $\bar c=2\gamma$, $\bar m=-\gamma^2$ put $P=1/\gamma$.
Section~\ref{rw:upper-degenerate-preparation} provides the exact
upper factor $f_U(v)=\mu_U+v^2$, its normally contracting radial
map, and the finite entry $v_e(p)$ with $v_e'<-c$.
The limiting entry is zero at $p=P$ and negative for $p>P$.
Throughout this subsection retain the full test
\begin{equation}\label{rw:infinity:quadratic-test}
 \mathcal P(\mu,v):\quad
 \mu>0\quad\hbox{or}\quad[\mu\le0\hbox{ and }v>\sqrt{-\mu}],
 \qquad v<v_R.
\end{equation}
This tests positivity on the entire interval to the fixed right exit:
a root between two positive endpoint values still blocks passage.
For $\bar B>0$, the horizontal power and the regular angular
connector give the complete incoming rate
\begin{equation}\label{rw:infinity:repeated-rate}
 J(v)=K_J\exp\left[-\frac{h_I}{\epsilon}
         \int_v^{v_R}\frac{\beta_U(s,0)}{f_U(s)}\,ds
                      +\frac{\Psi(v)}\epsilon\right],\qquad
 \frac{J''}{J'}=\frac{A_U(v)}{\epsilon f_U(v)}.
\end{equation}
Here $h_I=(1-\epsilon\bar B)/\bar B$ and
\[
 N=h_I\beta_U+f_U\Psi',\qquad
 A_U=N-\epsilon f_U'+\epsilon f_UN'/N>c.
\]
The relative integral of Section~4, composed with independent normal
inputs, bounds the prescribed mixed jets of $\Psi$, the regular
connector terms and the rate coefficients. Any section-independent
divergent contraction remains in $K_J$. The input-dependent normal
output enters only through its bounded absolute jets.

\begin{proposition}[Repeated upper root away from the joint center]
\label{rw:infinity:repeated-lower-loss}
Suppose $\bar B>0$, the finite lower polynomial loses positivity
at or before $p_*\ge P$, and
$(\bar B,\bar a)\ne(\gamma^2,-2\gamma)$.
The local closing count is at most three, and at most two unless
a simple lower root lies at $p_*=P$.
\end{proposition}
\begin{proof}
For a simple root at $P$, use its positive power input $t$.
The upper input $v(t)$ has $v'>c$ and tends to zero at the limiting
endpoint. Clearing the critical derivative gives
\[
 M=tC-f_U(v)H_o,\qquad
 C=A_U(v)v'+\epsilon f_U(v)v''/v'>c,
 \quad H_o=h_+-\epsilon+\epsilon tR_+(t)>c.
\]
Since $t,v,f_U$ are small, the bounded jets and $H_o'=O(\epsilon)$
give $M'>c$ on a fixed smaller interval. It has at most one zero.
If a simple or double lower root lies at $p_*>P$, its rate has
vanishing denominator while the upper entry factor has a positive
margin. After multiplying by the lower denominator, the critical
numerator is strictly negative. If $p_*=P$ and a double lower
root lies strictly before $P$, the lower entry factor instead has
a positive margin, and the upper rate dominates with a strictly
positive numerator. These cases give two closing zeros.

It remains to compare a double lower root $r=1/\sqrt{\bar B}$
and the repeated upper threshold $P$ when both precede $p_*$.
The entries then stay a fixed distance beyond their slow minima.
Subtract the divergent actions, which are constant in $p$, before
differentiating. The limiting logarithmic rates in this coordinate are
\[
 \epsilon\partial_p\log|D_{+,p}|\longrightarrow
                   -\frac p{\bar B(p-r)^2},\qquad
 \epsilon\partial_p\log|J_p|\longrightarrow
                   -\frac p{\bar B(p-P)^2}.
\]
For the first, the moving lower rate at its regular entry is
$-A_+'(p)$. For the second take the exact affine upper field with
$\bar m=-\gamma^2-\omega^2$. At fixed $p>P$ its return has
\[
 T=\pi/\omega-1/[\gamma^2(p-P)]+o(1),\qquad
 q=(p-P)e^{\gamma T}(1+o(1)).
\]
Thus $\partial_p\log q\to p/(p-P)^2$, and the incoming
horizontal rate multiplies it by $-1/\bar B$.
To pass from the affine calculation to the original maps, use the
positive margins of the reduced denominators on these entry intervals
and the common integrable majorants for the differentiated relative
normal integrals from Section~4. At the repeated limiting field
its exact $1/V$ normal coordinate makes the upper relative logarithm
zero, as established in Section~\ref{rw:upper-degenerate-preparation}.
The connector's nonconstant terms
are evaluated at a normal output tending to zero. These facts prove
convergence in the finitely many ordinary input jets used here,
without differentiating the divergent constant action.
Subtracting the two displayed rates gives
\[
 \frac{p(r-P)(2p-r-P)}
      {\bar B(p-r)^2(p-P)^2},
\]
whose sign is nonzero. Equality $r=P$ is exactly the excluded joint
center. The strict rate difference gives two closing zeros.
All the counts are on the intersection of the whole-path conditions
and the positive central-output interval, and include multiplicity.
\end{proof}

When $r=P$, neither rate dominates. Both signed gaps must then be
retained in the joint-center chart of Section~\ref{rw:joint-section}.
\label{rw:infinity:signed-horizontal}
We next cover $\bar B=0$, where replacing the horizontal factor
by a power with exponent $1/\bar B$ would not be legitimate.
Here a finite lower loss is simple and $\bar a<0$.
The exact phase of Proposition~\ref{rw:phase} has $F_I'>c$ near
zero. Set $u=F_I(x)-\bar B$ and divide time by the positive unit
$xF_I'(x)/u$. This produces
\[
 \dot u=\epsilon u(\bar B+u),\qquad
 \dot Z=-Z\widetilde\beta(u,Z),\qquad
 \widetilde\beta,\ \widetilde\beta+Z\widetilde\beta_Z>c.
\]
The unit extends through $u=0$ and both signs of $\bar B$.
The exact normal input supplies
\[
 \frac{D_I''}{D_I'}=\frac{A_I(u)}{\epsilon u(\bar B+u)},
 \qquad A_I>c,\qquad u>\max(0,-\bar B).
\]
The moving-entry formula bounds the ordinary and mixed jets of
$A_I$. This exact quadratic factor remains valid through
$\bar B=0$ without dividing by $\bar B$.

For a strict upper pair write the input as $u=u(t)$ and put
$\mathcal R=tu'/u=r_s+O_D(t^{\alpha_0})$,
$tu''/u'=r_s-1+E_1$, $E_1=O_D(t^{\alpha_0})$.
With the simple lower input $\delta+t$, set
\[
 Q=A_I(u)\mathcal R+\epsilon(\bar B+u)(r_s-1+E_1)>c,
 \quad H_o=h_+-\epsilon+\epsilon(\delta+t)R_+(\delta+t)>c.
\]
The full numerator is
\[
 M=(\delta+t)Q-t(\bar B+u)H_o.
\]
Its clearing factor is positive on
$t>0$, $\delta+t>0$, $u>\max(0,-\bar B)$.
At a zero $(\delta+t)/t\asymp\bar B+u$ is small. Since
$|tQ'|\le Ct^{\alpha_0}$, $tu'=u\mathcal R$ and
$H_o'=O(\epsilon)$, differentiation gives $M'=Q+o(1)>c/2$
there. The same-direction crossing argument gives at most one numerator
zero and three closing zeros, for both signs of nearby $\bar B$.

For a repeated upper pair let $u=u(v)$ include its regular angular
connector. Write
\[
 f(v)=\mu_U+v^2,\quad S(v)=\bar B+u(v),\quad
 L(v)=fu_v/u>c.
\]
The input-independent normal contraction is separated first.
Its relative formula and the bounded absolute jets of $u$ show that
$L$ and its needed ordinary jets are bounded. The exact incoming rate is
\[
 \mathcal B:=\epsilon fS\frac{J_{vv}}{J_v}
 =A_I(u)L+\epsilon S(L-f'+fL'/L)>c.
\]
In a positive finite lower power input $t$, the regular entry $v(t)$
has $v'>c$. Set
$C=\mathcal B(v)v'+\epsilon f(v)S(v)v''/v'$ and
$H_o=h_+-\epsilon+\epsilon tR_+(t)$. The numerator and derivative are
\[
 M=tC-f(v)S(v)H_o,\qquad
 M'=C+tC'-[2vv'S+uLv']H_o-fS H_o'.
\]
Use $fu_v=uL$ before estimating to cancel the apparent singularity
in the composition derivative. Fix the finite upper entry neighborhood
and its bounded jets first, then shrink $t$, $|\bar B|$ and the
normal output $u$. The result is $M'>c$ on the entire passing
interval, hence three closing zeros. Both
\eqref{rw:infinity:quadratic-test} and $u>\max(0,-\bar B)$ remain
imposed. This also covers finite entries strictly beyond $P$.

\subsection{The inverse upper map at the primary weak boundary}
\label{rw:infinity:weak-section}

At the primary weak boundary, the upper return is best described by
its inverse: the absolute output becomes flat, while the scalar factor
still records the passing side. Fix $\bar m=0$, $\bar c=c_0>0$
and $P=1/c_0$.
Use the actual weak parameter and normal form of
\eqref{rw:upper:weak-parameters} and
Section~\ref{rw:upper-degenerate-preparation}; the actual weak face is
$\bar m+\epsilon\bar B(\bar c+\nu)^2=0$.
In reversed projective time the exact scalar factor is
$f(u)=u(\gamma+u)$ and the fast direction contracts uniformly.
The complete inverse upper displacement, including its regular finite
connection, is
\begin{equation}\label{rw:infinity:weak-upper}
 p=\widehat p-t(u),\qquad
 \frac{t'}t=\frac Nf,\quad \frac{t''}{t'}=\frac{A_U}f,
 \qquad N,A_U>c,\quad u>\max(0,-\gamma).
\end{equation}
Here $N=\beta_U+f\widetilde j'$, $A_U=N-f'+fN'/N$.
Their required ordinary and mixed jets are bounded. The absolute output
$t$ and its selected derivatives tend uniformly to zero when the
input and $\gamma$ neighborhoods shrink. At a lost exit, extend the
absolute output by zero and the relative coefficients by the complete auxiliary flow. Coefficient division uses
these extensions rather than an inverse of the flat map.

The incoming horizontal entry $x=\chi(u)$ is the regular connector
of Proposition~\ref{rw:center-coordinates}. Thus
\[
 \chi(0)=0,\quad\chi'>c,\quad
 L=u\chi'/\chi,\quad L(0)=1,\quad R_\chi=\chi''/\chi'
\]
have bounded required jets. Keep the exact raw factor of
Proposition~\ref{rw:phase}, with no division by $\bar B$ or
$\bar a$. Its complete rate is
\begin{equation}\label{rw:infinity:raw-horizontal}
 \frac{D_I''(x)}{D_I'(x)}
       =\frac{A_I(x)}{\epsilon xF_I(x)},\qquad
 A_I=1+O_{C^N}(\epsilon),\quad A_I(0)=1-2\epsilon\bar B.
\end{equation}
The whole condition is $F_I(s)>0$ for $\chi(u)\le s\le x_R$.
Since $F_I''>1$ and $F_I(x_R)>c$, its retained side is the interval
to the right of its last root, or the whole input interval if there
is no root. This remains valid at the cubic boundary
$\bar B=\bar a=0$, where checking the entry value alone would
not exclude a later root.

If an incoming finite double root is separated from the horizontal
input, put fixed sections on both sides of it. Its whole-passage
condition is a positive minimum, independent of $u$. The relative
normal integral separates its possibly divergent contraction into
an input-independent linear factor. Estimate the horizontal output first
and then compose the bounded relative unit. This preserves \eqref{rw:infinity:raw-horizontal}
and its required ordinary input jets. The divergent constant is not
differentiated in coefficients. If the whole-passage condition fails,
the corresponding input set is empty.

\begin{proposition}[A regular outgoing arc at the weak boundary]
\label{rw:infinity:weak-regular}
At $\bar m=0$, $\bar c>0$, $p_*=P$, $q_*=\infty$, a regular
finite outgoing lower arc has local closing count three. An incoming
finite double root, either sign of $\ell$, and the reflected case
are included.
\end{proposition}
\begin{proof}
The regular outgoing deviation $D_o(p)$ has
$R_o=-\epsilon D_o''/D_o'>c$ and bounded needed ordinary input
jets, by Lemma~\ref{rw:finite:regular-lower}. Put
$I(u)=D_o(\widehat p-t(u))$ and $J(u)=D_I(\chi(u))$.
Both have positive derivatives. For $K=\epsilon\log(J'/I')$,
direct differentiation gives the exact numerator
\begin{equation}\label{rw:infinity:weak-regular-numerator}
\begin{aligned}
 M&=uF_I(\chi)(\gamma+u)K'\\
  &=A_I(\chi)L(\gamma+u)-\epsilon F_I(\chi)A_U
    +\epsilon uF_I(\chi)(\gamma+u)R_\chi
    -F_I(\chi)R_o(\widehat p-t)Nt.
\end{aligned}
\end{equation}
Apply Lemma~\ref{rw:weak:rate-calculus} with
$h=F_I(\chi)$, $A=A_I L$, $R=F_I(\chi)R_\chi$ and
$E=F_I(\chi)R_oN$. Their required $C^1$ bounds follow from the
prepared rates and regular outgoing arc; $A=1+O(u+\epsilon)$ and
$\|t\|_{C^1}\to0$. It gives $M'>1/2$ and three closing zeros by
Lemma~\ref{rw:zero-calculus}, on the intersection of the whole upper,
horizontal and positive-central passing conditions. A separated double
passage imposes its additional coefficient condition. All nearby
counted cycles with $q\to\infty$ reach this horizontal crosscut
with $u\to0$ and follow the unique inverse upper transition.
\end{proof}

\subsection{Finite lower loss at the weak boundary}

The weak upper output may also meet a lower loss. First keep
$\bar B_0>0$. The necessary outgoing alternatives are
a first simple root at $P$, with
$\bar a=-\bar c-\bar B/\bar c$ and $0<\bar B<\bar c^2$;
a double root at $P$, with $\bar a=-2\bar c$,
$\bar B=\bar c^2$; or a double root strictly before $P$.
Here $\bar a<0$, so the incoming half-line is regular.
For a root at $P$, choose a decreasing lower slow input $v_L(p)$,
put $\delta_L=v_L(\widehat p)$ and replace $t$ by
$v_L(\widehat p-t)-v_L(\widehat p)$. This is a positive regular
change of the absolute upper displacement; its rates still have the
form \eqref{rw:infinity:weak-upper} with bounded positive $N,A_U$.
The lower rate is
\[
 \frac{D_L''(v)}{D_L'(v)}=\frac{A_L(v)}{\epsilon f_L(v)},\qquad
 f_L(v)=v\quad\hbox{or}\quad f_L(v)=\mu_L+v^2,
 \qquad A_L>c.
\]
For a double root impose positivity of $f_L$ throughout its entire
slow path. Both inputs now have positive derivatives:
$I(u)=D_L(\delta_L+t(u))$ and $J(u)=D_I(\chi(u))$.

Their critical numerator is
\begin{equation}\label{rw:infinity:weak-loss-numerator}
\begin{aligned}
 M&=f_L(\delta_L+t)C-Dt,\\
 C&=A_I(\chi)L(\gamma+u)-\epsilon F_I(\chi)A_U
             +\epsilon uF_I(\chi)(\gamma+u)R_\chi,\\
 D&=F_I(\chi)A_L(\delta_L+t)N.
\end{aligned}
\end{equation}
Its clearing factor is $uF_I(\chi)(\gamma+u)f_L(\delta_L+t)>0$.
The upper inverse curvature contributes $-\epsilon F_I A_U$ to
$C$; it remains in all subsequent rate comparisons.

\begin{proposition}[Weak boundary with positive $\bar B$ and lower loss]
\label{rw:infinity:weak-loss}
At a root at $P$ the unanchored local bounds are $Z(K')\le2$,
$Z(K)\le3$ and $Z(F)\le4$, including multiplicity.
A double root strictly before $P$ has closing count three.
\end{proposition}
\begin{proof}
The coefficients in \eqref{rw:infinity:weak-loss-numerator} have
bounded $C^2$ input norms, $C'>1/2$, and arbitrarily small $|C|$.
Moreover $D>c$ because $\bar B_0>0$.
Write $f_L(\delta_L+t)=A+rt$, with
$(A,r)=(\delta_L,1)$ or
$(\mu_L+\delta_L^2,2\delta_L+t)$, and put $E=D-rC>c$.
Then $M=AC-tE$. At any zero the original expression
$f_L C=Dt$ forces $C>0$, hence also $A>0$.
For $A>0$ define $R_0=C/(tE)$ on the entire passing interval.
Its derivative has sign opposite to
\[
 V=NCE+f(CE'-EC'),\qquad
 R_0'=-V/(ftE^2).
\]
One differentiation gives
\[
 V'=NC'E+N'CE+NCE'+f'(CE'-EC')+f(CE''-EC'')>0.
\]
The first term has a positive margin and the others are
$O(|C|+|f'|+|f|)$, so $V'>c$ on a common neighborhood.
The argument uses only small $C$, $C'>c$, $E>c$ and their bounded
$C^2$ norms; this is the quotient estimate used again at the weak center.
Every critical point of $C/(tE)$ is a strict maximum, since
$R_0''=-V'/(ftE^2)<0$ there. Its strict superlevel sets are intervals
and each level has at most two intersections with multiplicity.
At $C\le0$, $V=C(NE+fE')-fEC'<0$, so every nonpositive-level
crossing is strictly increasing. In particular the level $1/A$
proves the first count.
If the double root precedes $P$, use fixed sections around it and
separate its input-independent contraction. The outgoing entry rate
is regular with positive $R_o$; the proof of
Proposition~\ref{rw:infinity:weak-regular} then applies unchanged.
\end{proof}

\subsection{Center division through \texorpdfstring{$\bar B=\bar m=0$}{B=m=0}}
\label{rw:infinity:weak-center-section}

At $\bar B=\bar m=0$, the preceding positive lower bound for $D$
is lost. Near the remaining threshold $\bar a=-\bar c<0$, divide
by the center generators in the weak coefficient box of
Proposition~\ref{rw:center-coordinates}:
\[
 (\bar B,\bar c,d_s,\varsigma,\nu),\qquad
 d_s=P_L-\widehat p,\quad\varsigma=\bar B+\bar m.
\]
It includes both signs of $\bar B$ and $\bar m$ and has center
axis $d_s=\varsigma=\nu=0$. The lower coordinate is the unnormalized
deficit $v=P_L-p$. Its simple-root map has the form
\[
 D_L(v)=K_Lv^\lambda e^{\Psi_L(v)/\epsilon},\qquad
 \frac{D_L''}{D_L'}=\frac{A_L(v)}{\epsilon v},\quad
 A_L(0)=\epsilon(\lambda-1),\quad\epsilon\lambda>c.
\]
The scaled unit $\Psi_L$ has bounded mixed jets, since a regular
power coordinate $v k(v)$ adds $\epsilon\lambda\log k(v)$.
The unscaled logarithm need not be bounded. Since $\bar a<0$
has a fixed margin, the raw horizontal factor is increasing here.

The actual interval is precisely
\begin{equation}\label{rw:infinity:weak-center-domain}
 u>0,\qquad\gamma+u>0,\qquad F_I(\chi(u))>0,\qquad d_s+t(u)>0,
\end{equation}
together with positive central outputs. All four displayed conditions
are one-sided monotone restrictions. Put $v=d_s+t(u)$ and reuse $C,D$
from \eqref{rw:infinity:weak-loss-numerator}, with $f_L(v)=v$.
The exact numerator is
\begin{equation}\label{rw:infinity:weak-center-numerator}
 M=d_s C-tE,\qquad E=D-C,\qquad
 M=uF_I(\chi)(\gamma+u)(d_s+t)K'.
\end{equation}
We have $C'=1+O(u+|\gamma|+\epsilon)$ and a small bounded $C$.
The following center division replaces the missing positive lower
bound for $D$ across $\bar B=0$.

\begin{proposition}[The weak-center local bound]
\label{rw:infinity:weak-center}
On a common neighborhood of this threshold, every noncenter fibre
has at most four isolated closing zeros, counting multiplicity,
on its actual passing interval. On a fibre with a closing zero,
\begin{equation}\label{rw:infinity:weak-relative-forcing}
 |\nu|\le\eta(\epsilon)(|d_s|+|\varsigma|),\qquad\eta(\epsilon)\to0.
\end{equation}
If the two generators vanish, the field is a center and contributes
no isolated zero.
\end{proposition}
\begin{proof}
Reserve sufficient primitive orders for $C_u^2C_{\mathbf b}^2$
of the multiplied rates, including the two input derivatives spent
in each logarithmic derivative. For example total order ten in the
relative maps and connections suffices. The lower large power, on
a fixed small deficit interval, and the horizontal normal contraction
followed by the fixed central core give
\[
 \|D_L\|_{C_v^4C_{\mathbf b}^2}
       +\|D_I\|_{C_x^4C_{\mathbf b}^2}\le Ce^{-c/\epsilon}
\]
after reducing $c$ to absorb the finite inverse-scale powers.
For the first term choose the deficit radius so its logarithm
dominates the bounded scaled amplitude. For the second keep a fixed
positive-time final segment of the horizontal/core passage. Its
normal contraction supplies a positive action independent of all
weak gaps. Finite parameter variations are bounded by the normally
dominated relative construction; the same exponential absorbs their
polynomial losses. Both absolute maps have the zero extensions of
Section~4. Their rate coefficients have bounded higher mixed
extensions on the larger auxiliary coefficient domain.

For $A=|d_s|+|\varsigma|+|\nu|$, Proposition~\ref{rw:domains:weak}
supplies the single real center
$(\bar B+K_*A,\bar c,0,0,0)$ through every original passing input.
Its complete deviations agree there, and therefore $E=0$.
Comparison of the absolute maps on that common domain yields
$|I-J|\le CAe^{-c/\epsilon}$. With the fixed central gap
$\rho\nu H$, $H>h_*$, Lemma~\ref{rw:centers:relative-forcing}
proves \eqref{rw:infinity:weak-relative-forcing}.

The same center supplies the rate identity needed for division.
Each point of the original interval lies in a strictly passing open
center neighborhood, so both $E$ and its passive $\bar B$
derivative vanish there. Taylor expansion gives
\[
 E=d_s E_d^c+\varsigma E_\varsigma^c+\nu E_\nu^c+R,\qquad
 \|R\|_{C_u^2}\le CA^2.
\]
For a fixed nonzero transverse vector $g=(d_s,\varsigma,\nu)$
set $e_i=E_i^c+g_iR/|g|^2$.
These provide an exact division $E=d_s e_d+\varsigma e_\varsigma+\nu e_\nu$
with bounded $C_u^2$ norms and error $O_{C_u^2}(A)$ from the center
columns. Only the fixed-fibre division is used; smooth dependence of
this choice at $g=0$ is unnecessary.

To determine the sign of the $\varsigma$ column, use the exact
endpoint identity. At $u=t=d_s=0$, the actual coefficients satisfy
\[
 N(0)=\Delta,\quad A_U(0)=\Delta-\gamma,\quad
 F_I(0)=\bar B,\quad A_I(0)=1-2\epsilon\bar B,
\]
so
\begin{equation}\label{rw:infinity:weak-endpoint-E}
 E(0)=\epsilon\bar B\lambda\Delta-(1-\epsilon\bar B)\gamma.
\end{equation}
At $\bar B=0$ this is $-\gamma$ for every $d_s$.
Since $\bar m=\varsigma-\bar B$, at the zero center
$\partial_\varsigma E(0)=1/\bar c>0$ and
$\partial_{d_s}E(0)=0$, at each small positive scale.
The bounded mixed extension then gives $e_\varsigma>c>0$
on a fixed smaller domain.

Absorb the $d_s$ column into $\overline C=C-te_d$.
It is small, bounded in $C^2$, and $\overline C'>1/2$.
The numerator becomes
\[
 M=d_s\overline C-t(\varsigma e_\varsigma+\nu e_\nu).
\]
On a fibre with a closing zero put $I_0=|d_s|+|\varsigma|$.
If $|\varsigma|\ge\sqrt\eta I_0$, then
$e=e_\varsigma+(\nu/\varsigma)e_\nu>c/2$.
For $d_s=0$ there is no numerator zero. Otherwise apply the quotient
calculation of
Proposition~\ref{rw:infinity:weak-loss} with
$(C,E)=(\overline C,e)$. It makes
$R_0=\overline C/(te)$ unimodal and gives at most two intersections
with the level $\varsigma/d_s$. Its hypotheses hold because
$\overline C$ is small, $\overline C'>1/2$, $e>c/2$ and both
have the required bounded $C^2$ norms.
In the complementary case $|d_s|\ge(1-\sqrt\eta)I_0>0$, and
\[
 M/d_s=\overline C-t[(\varsigma/d_s)e_\varsigma
                                      +(\nu/d_s)e_\nu]
\]
is strictly increasing. Hence $M$ has at most two zeros in either
case. Its positive multiplier and two Rolle steps give four closing
zeros. The auxiliary extensions enter only the coefficient estimates.
The zeros counted by Rolle's theorem lie in
\eqref{rw:infinity:weak-center-domain}.
\end{proof}

\begin{lemma}[Weak-chart entrance data]
\label{rw:infinity:weak-entry-data}
The weak charts above supply the following data independently of any
preceding zero count. Outward motion is increasing $p$ and decreasing
$u$, and $\operatorname{sign}K_p=-\operatorname{sign}M$.
In the positive-$\bar B$ lower-loss chart, a numerator zero requires
$A>0$ and $C>0$, and the strict superlevel sets of $C/(tE)$ are
intervals. The limiting regular critical function is positive on
its whole regular interval; this gives a positive actual margin on
any retained compact regular overlap.

In the weak-center chart, on each noncenter fixed fibre satisfying
\eqref{rw:infinity:weak-relative-forcing}, on the same actual passing
interval, the coefficient division in
Proposition~\ref{rw:infinity:weak-center} gives
\[
 M=d_s\overline C-t(\varsigma e_\varsigma+\nu e_\nu),
 \qquad \overline C=C-te_d,\qquad
 e_\varsigma>c>0,\quad\overline C'>1/2.
\]
The divided coefficients have bounded $C_u^2$ norms and
$\overline C$ is small on the same common neighborhood; the division
is chosen on each fixed nonzero transverse vector.
Put $I_0=|d_s|+|\varsigma|$. If
$|\varsigma|<\sqrt\eta I_0$, then $d_s\ne0$ and $M/d_s$ is
strictly increasing. Otherwise $\varsigma\ne0$ and
$e=e_\varsigma+(\nu/\varsigma)e_\nu>c/2$.
For $d_s\ne0$ in this alternative, $R_0=\overline C/(te)$
increases and then decreases, either portion possibly absent.
Its strict superlevel sets are intervals, and every nonpositive-level
crossing is strictly increasing. Each level has at most two
intersections, counted with multiplicity; in particular $Z(M)\le2$.
The identity
\[
 M=d_s t e(R_0-\varsigma/d_s)
\]
records which of those superlevel sets contains an entrance with a
prescribed sign of $M/d_s$. The norms of $(d_s,\varsigma)$ and
$(\tau,\sigma_1)$, $\sigma_1=\bar c(\bar B+\bar m)$, are
comparable under the relative forcing condition. These facts will
be used to anchor the terminal interval in Section~12.
\end{lemma}
\begin{proof}
The first statement follows from $p=\widehat p-t(u)$, $t'>0$,
and the positive multipliers defining $M$.
The quotient calculation in Proposition~\ref{rw:infinity:weak-loss}
gives the superlevel, multiplicity and nonpositive-level assertions
for both $(C,E)$ and $(\overline C,e)$.
Norm comparison is \eqref{rw:centers:single-ideal-comparison},
with its $|\nu|$ term absorbed using
\eqref{rw:infinity:weak-relative-forcing}.

For the remaining regular sign at $\bar B>0$, the threshold
coefficients have $\tau<0$ and $\sigma_1=\bar c\bar B>0$.
At $\bar m=0$ the exact limiting upper return is parametrized by
\[
 p=P-\frac{T}{e^{\bar cT}-1},\qquad q=T-p.
\]
For large $T$, $A_-(q)=O(\log T)$. At a simple outgoing root
$A_+(p)\sim\bar cT/(\bar c^2-\bar B)$; at a double root
$A_+(p)\sim e^{\bar cT}/(\bar c^3T)$.
These follow by integrating the respective simple or double pole
of $p/F_+(p)$ and inserting $P-p\sim Te^{-\bar cT}$.
Thus $K^0=A_+(p)-A_-(q)>0$ near $P$. It is positive near zero
as well, by \eqref{rw:finite:small-columns} and $\tau<0$.
Lemma~\ref{rw:regular-columns} allows at most one simple zero,
so there is none between these two ends. The finite-scale regular
convergence proves the positive margin on each fixed compact overlap.
This is the entry sign used in Section~12 to count the regular and
weak parts together; their separate local bounds are not added.
\end{proof}

\begin{proposition}[Local exhaustion for one infinite endpoint]
\label{rw:infinity:one-infinite}
Every necessary limiting configuration with exactly one finite and one
infinite ray endpoint, except the joint center reserved for
Section~\ref{rw:joint-section}, has a neighborhood with a common
small-scale bound of four isolated counted cycles, including multiplicity.
The neighborhoods are open in the original normalized coefficients
and include their signed slow gaps and both line-breaking signs.
\end{proposition}
\begin{proof}
Lemma~\ref{rw:finite:upper-return} and its strict barrier give the
three upper alternatives stated above. For $\bar m<0$, apply
Proposition~\ref{rw:infinity:negative-regular} to a regular finite
lower arc. A lower loss has $\bar a<0$, $\bar B\ge0$ and is a
first simple root at the endpoint or a double root at or before it.
Propositions~\ref{rw:infinity:strict-lower-loss} and
\ref{rw:infinity:repeated-lower-loss}, together with the signed
horizontal construction at $\bar B=0$, cover these losses except
for the reserved joint center.

At $\bar m=0$, regular finite data are covered by
Proposition~\ref{rw:infinity:weak-regular}. Positive-$\bar B$
losses are Proposition~\ref{rw:infinity:weak-loss}. At
$\bar B=0$ the two lower nonnegativity conditions give
$-\bar c\le\bar a\le0$: strict inequality at the left endpoint
is regular, and equality is the weak-center chart of
Proposition~\ref{rw:infinity:weak-center}. The raw horizontal
construction includes $\bar a=0$ and its cubic factor.
For an outgoing infinite endpoint reflect the field. The lower
restrictions of Lemma~\ref{rw:finite:action-order} require a loss
at or before its new finite outgoing endpoint; all those proofs
allowed either sign of the reflected $\ell$.

An empty whole-path domain contributes no cycle. On a nonempty
physical interval the original maps are analytic transverse
finite-time maps. The finite-smooth auxiliary coordinates and positive
multipliers used in Rolle's theorem preserve their zero multiplicities.
Reflection preserves the same return-germ multiplicity.
The neighborhood and scale cutoff in each argument are fixed near
one one-infinite limit. When the finite endpoint also tends to
infinity, the next subsection supplies the required neighborhoods.
\end{proof}
